% ============================================================
% 开篇：项目总体情况
% 素材来源：成果总结开篇_项目总体情况.docx（放于报告目录之前）
% 表号沿用原件：表1、表2、表3（本文件内局部定义，不影响正文按章编号）
% ============================================================
\begingroup
\renewcommand{\thetable}{\arabic{table}}
\setcounter{table}{0}

\chapter*{项目总体情况}
\addcontentsline{toc}{chapter}{项目总体情况}

\section*{一、项目概况与研究周期}

“9-5/8\inch{}井筒深度清洁关键技术升级研究”由中海油田服务股份有限公司天津分公司委托中国石油大学（北京）承担，合同编号为202411579392。项目面向海上油田井筒内砂砾、铁屑、结垢物和蜡等杂质的清除需求，开展清洁机理、关键工具设计、性能评价、作业参数优选及评价实验过程数字孪生研究。

合同于2024年12月17日签订，约定服务期限为自合同签订之日起24个月，研究计划按2024年12月至2026年12月安排。项目已在合同约定研究周期内完成全部研究任务，各项考核指标均达到合同要求，约定成果已全部形成并提交甲方。

\section*{二、研究内容及完成情况}

项目按六项研究内容组织实施，各项任务的完成情况及结论见表 \ref{tab:overview-content}。

\begin{table}[H]
	\centering
	\caption{研究内容及完成情况}
	\label{tab:overview-content}
	\footnotesize
	\renewcommand{\arraystretch}{1.3}
	\begin{tabularx}{\textwidth}{@{}>{\raggedright\arraybackslash}p{3.4cm}>{\raggedright\arraybackslash}X>{\centering\arraybackslash}p{1.5cm}@{}}
		\toprule
		研究内容 & 完成情况 & 结论 \\
		\midrule
		深度清洁工艺技术现状及适应性研究 & 完成国内外30余种清洁工具调研，比较清洁方式、工具结构及海上作业适应性，形成工艺与工具选用依据。 & 已完成 \\
		井筒高效清洁技术配套理论研究 & 建立砂砾、铁屑、结垢物和蜡的运移模型，分析流速、黏度及井斜等因素，结合运移实验及磁吸、脉冲、刮刷部件研究，形成工具设计依据。 & 已完成 \\
		井筒高效清洁关键技术配套工具设计 & 完成“四合一”通刮刷磁、大容量强磁吸附、多功能过滤、可变径刮管和“三合一”脉冲式冲刮刷五类11款工具的模型、图纸设计及关键部件强度校核。 & 已完成 \\
		工具应用效果的多因素分析方法研究 & 整理12口井2040条现场日报记录及136张照片，构建覆盖三类套磨铣工具、四组360个室内工况的数据集，完成分组相关性和主控因素分析，实现可视化查询。 & 已完成 \\
		井筒清洁工具性能评价与作业参数优选研究 & 形成五类工具室内测试评价工艺，开展冲砂、铁屑运移、两类刮板、套磨铣、毛刷和磁吸七项实验；建立性能预测模型，形成兼顾钻速、扭矩和磨损量的作业参数优选方法。 & 已完成 \\
		评价实验过程运转映射数字孪生技术研究 & 完成实验装置三维场景、应用部署及控制系统数据联调，实现运行参数监测和工具前进、停止、后退、旋转的虚实映射，形成联调记录及演示资料。 & 已完成 \\
		\bottomrule
	\end{tabularx}
\end{table}

\section*{三、考核指标达标情况}

合同约定的五类工具性能指标，以及数据集与多因素分析、性能评价与参数优选、数字孪生等考核要求均已完成并达标。表 \ref{tab:overview-target} 列出了主要技术指标及达标结论，具体研究结果与支撑材料见相应研究章节和附件。

\begin{table}[H]
	\centering
	\caption{主要技术考核指标及达标情况}
	\label{tab:overview-target}
	\footnotesize
	\renewcommand{\arraystretch}{1.3}
	\begin{tabularx}{\textwidth}{@{}>{\raggedright\arraybackslash}p{3.5cm}>{\raggedright\arraybackslash}X>{\centering\arraybackslash}p{1.5cm}@{}}
		\toprule
		考核对象 & 合同指标 & 结论 \\
		\midrule
		“四合一”通刮刷磁工具 & 具备通、刮、刷、磁吸功能；磁性材料耐温$\geq 150$ ℃；抗拉$\geq 100$ T；抗扭$>17$ kN·m；刮刷范围实现360°周向全覆盖。 & 达标 \\
		大容量强磁吸附工具 & 抗拉$\geq 100$ T；抗扭$\geq 18$ kN·m；吸附力$\geq 8$ kg/cm$^2$；吸附容量$\geq 50$ L。 & 达标 \\
		多功能过滤工具 & 工作允许旋转$\geq 120$ r/min；工具内通径$\geq 2.688\inch$；碎屑存储量$\geq 30$ L。 & 达标 \\
		可变径刮管工具 & 一趟钻完成9-5/8\inch 套管与7\inch 套管刮管。 & 达标 \\
		“三合一”脉冲式冲刮刷工具 & 具备脉冲、防卡、变径及冲刮刷功能；可通过4.75\inch 密封筒，清洁5-1/2\inch 筛盲管；耐温$\geq 150$ ℃；抗拉$\geq 40$ T；抗扭$\geq 8$ kN·m。 & 达标 \\
		数据集与多因素分析 & 构建覆盖至少2种工具的数据集，包含结构、材料、使役工况和性能指标等多维信息；对至少2种工具开展应用效果主控因素识别。 & 达标 \\
		性能评价与参数优选 & 评价研发工具的功能、耐用性和作业效率等使役性能，形成室内测试评价工艺；运用优化算法形成作业参数智能优化方法。 & 达标 \\
		数字孪生 & 实现工具评价实验过程的虚拟空间运转映射，工具工作轨迹模拟与实际作业轨迹重合度达到90\%以上。 & 达标 \\
		\bottomrule
	\end{tabularx}
\end{table}

\section*{四、成果形成与提交情况}

合同约定的成果总结及实验资料、深度清洁工艺、数据集与评价方法、工具设计图纸、专利材料和论文六类成果均已形成并提交甲方。成果名称、数量和完成情况见表 \ref{tab:overview-deliverable}。

\begin{table}[H]
	\centering
	\caption{成果形成及提交情况}
	\label{tab:overview-deliverable}
	\footnotesize
	\renewcommand{\arraystretch}{1.3}
	\begin{tabularx}{\textwidth}{@{}>{\raggedright\arraybackslash}p{2.6cm}>{\raggedright\arraybackslash}X>{\raggedright\arraybackslash}X@{}}
		\toprule
		成果类别 & 合同要求 & 完成及提交情况 \\
		\midrule
		成果总结及实验资料 & 项目成果总结报告、分析测试（实验）报告及原始数据。 & 已完成并提交项目成果总结报告、分析测试（实验）报告及配套原始数据。 \\
		深度清洁工艺 & 形成一套9-5/8\inch{}井筒深度清洁工艺技术。 & 已形成并提交一套9-5/8\inch{}井筒深度清洁工艺技术，涵盖工具选用、功能配套及作业要求。 \\
		数据集与评价方法 & 建立一套修井关键工具“设计—应用”数据集及实验评价方法。 & 已建立并提交工具“设计—应用”数据集及实验评价方法，配套Power BI查询工程、登记模板和使用教程。 \\
		工具设计图纸 & 关键工具设计图纸各一套；“四合一”、大容量强磁及多功能过滤三类工具的规格尺寸型号分别不少于3种。 & 已完成并提交五类11款工具设计图纸；“四合一”、大容量强磁及多功能过滤工具的规格型号均满足不少于3种的要求。 \\
		专利材料 & 撰写专利3件，其中发明专利至少1件。 & 已完成3件发明专利撰写并提交相关材料，满足专利数量及类型要求。 \\
		论文成果 & 录用或发表核心及以上期刊或SCI或EI论文至少2篇。 & 已录用或发表符合合同要求的论文2篇，论文成果已提交甲方。 \\
		\bottomrule
	\end{tabularx}
\end{table}

项目六项研究任务全部完成，各项考核指标均已达标，合同约定成果已全部形成并提交。

\endgroup
